% !TEX TS-program = xelatex
% !TEX encoding = UTF-8 Unicode
% !Mode:: "TeX:UTF-8"

\documentclass{resume}
\usepackage{zh_CN-Adobefonts_external} % Simplified Chinese Support using external fonts (./fonts/zh_CN-Adobe/)
%\usepackage{zh_CN-Adobefonts_internal} % Simplified Chinese Support using system fonts
\usepackage{linespacing_fix} % disable extra space before next section
\usepackage{cite}

\begin{document}
\pagenumbering{gobble} % suppress displaying page number

\name{华祎恺}

% {E-mail}{mobilephone}{homepage}
% be careful of _ in emaill address
\contactInfo{(+86) 185-1612-0602}{huayikai.david@gmail.com}{}{https://github.com/DavidHua04}

% \\section{个人总结}
% 内容。

% \\section{\faGraduationCap\ 教育背景}
\section{教育背景}
\datedsubsection{\textbf{英属哥伦比亚大学},计算机科学,\textit{理学学士}}{2024.8 - 2027.5}
\ \textbf{绩点4.33 / 4.33 (91.5/100)}, 理学院国际生奖学金,院长学者
\datedsubsection{\textbf{香港中文大学（深圳）},计算机工程,\textit{工学学士}}{2023.6 - 2024.8}
\ \textbf{绩点3.7 / 4.0}

% \\section{\faCogs\ IT 技能}
\section{技术能力}
% increase linespacing [parsep=0.5ex]
\begin{itemize}[parsep=0.2ex]
  \item \textbf{编程语言}: Python (Pandas, Scikit-learn, TensorFlow, PyTorch, NumPy), Java (JavaFX), C, C++, JavaScript/TypeScript, SQL, R
  \item \textbf{测试}: JUnit, unittest, pytest, SuperTest, chai, Playwright
  \item \textbf{数据库}: MySQL, Oracle, PostgreSQL
  \item \textbf{Web开发}: RESTful服务, Node.js, Express, React, HTML, CSS, JavaScript/TypeScript
  \item \textbf{工具}: Docker, Git, Redis, Kubernetes, GitHub, Linux基础环境, CI/CD流水线
  \item \textbf{UI/UX设计}: 原型设计 (Figma), 访谈/观察/问卷设计, 主题分析, 可用性测试, 启发式评估
  \item \textbf{机器学习与人工智能}: 经典机器学习模型 (KNN, 岭回归, 随机森林, 集成方法等), 推荐系统, 时间序列, 模型评估与部署
\end{itemize}

\section{实习经历}
\datedsubsection{\textbf{吉赫兹科技有限公司 | 质量保证实习生}, 上海}{2025.5 - 2025.6}
\begin{itemize}
  \item 实习初期使用 \textbf{Figma} 参与 UI 设计迭代，根据客户需求变化调整界面布局。
  \item 独立采用 \textbf{敏捷}方法对某医院移动应用进行全流程端到端测试。
  \item 设计并维护100余条结构化测试用例，共发现21个缺陷及4个可用性问题，其中包括4个内部 QA 工程师遗漏的问题。
  \item 与 QA 工程师及开发人员协作，使用内部追踪工具对问题进行分类、上报及验收。
\end{itemize}

\datedsubsection{\textbf{香港中文大学（深圳）| 本科科研助理}, 深圳}{2024.1 - 2024.7}
\begin{itemize}
  \item 开发并实现网络爬虫，自动从企业官网获取 ESG 报告，显著提升数据采集效率。
  \item 将 PDF 文档转换为文本格式，提高数据准确性与可靠性。
  \item 调用 OpenAI API 从文本中提取关键信息，优化数据处理工作流。
  \item 设计精细化提示词，将幻觉或不准确数据的发生率降低 4\%，增强研究成果的完整性。
\end{itemize}

\section{项目经历}
\begin{itemize}[parsep=0.2ex]
  \item \textbf{基于LLM的ASCII艺术识别增强系统} (\textit{https://github.com/DavidHua04/ASCII-Art-Identifier})
  \begin{itemize}
    \item 设计轻量级内容审核工具，整合 LLM 与 VLM 模型以检测攻击性 ASCII 艺术内容。
    \item 通过图像重路由方案，将识别准确率从 12\% 提升至 81\%，幻觉率从 88\% 降至 14\%。
    \item 使用 Python 实现可复现实验及统计验证（t 检验、自助法）。工具：Python, OpenRouter
  \end{itemize}
  \item \textbf{UBC课程可视化规划工具} (\textit{https://github.com/DavidHua04/UBC-Courses-Visualization})
  \begin{itemize}
    \item 采用\textbf{测试驱动开发}方法，使用 \textbf{Playwright} 与 \textbf{chai} 开发课程可视化与规划工具。
    \item 跨平台部署（专属网站及PC应用），进行可用性测试以改善用户体验。工具：React, Tailwind.css, TypeScript, Express, Figma
  \end{itemize}
  \item \textbf{信用卡违约预测} (\textit{https://github.com/DavidHua04/CreditCardDefaultPost.github.io})
  \begin{itemize}
    \item 在3万条信用卡记录上构建并对比7种分类模型（逻辑回归、随机森林、XGBoost、LightGBM等），预测违约风险；通过 SHAP 分析确认还款延迟为主要预测因子。
    \item 经5折交叉验证与 RandomizedSearchCV 调优，选出 LightGBM 模型，测试集 ROC-AUC 达 0.78。工具：Python (Pandas, NumPy, scikit-learn, LightGBM, XGBoost, SHAP, Matplotlib)
  \end{itemize}
  \item \textbf{UBC Minecraft玩家参与度分析} (\textit{https://github.com/DavidHua04/Minecraft-Player-Engagement-Analysis})
  \begin{itemize}
    \item 利用玩家及会话数据分析 UBC 自建 Minecraft 服务器上的玩家行为，识别贡献数据最多的玩家类型，为后续研究招募提供参考。工具：Python (NumPy, Pandas, Matplotlib, Seaborn, SciPy)
  \end{itemize}
\end{itemize}

\section{其他经历}
\begin{itemize}[parsep=0.2ex]
  \item \textbf{英属哥伦比亚大学 | 本科助教}, 温哥华 \hfill 2026.1 - 2026.4
  \begin{itemize}
    \item 与另外两名助教共同主持每周实验课，协作讲解复杂编程概念并现场调试代码。
    \item 通过 Piazza 和答疑时间提供可扩展的技术支持，对常见问题创建可复用解答。
    \item 与教学团队协调统一评分标准，确保多个实验班次评估结果的一致性。
  \end{itemize}
  \item \textbf{香港中文大学（深圳）IEEE学生分会 | 财务官}, 深圳 \hfill 2023.8 - 2024.7
  \begin{itemize}
    \item 为200余名成员的组织设计并实施财务追踪系统，优化报销流程并确保合规。
    \item 主导跨部门协作，推动 IEEE 会员费报销自动化，协调学生分会与理工学院，提升会员参与率。
  \end{itemize}
\end{itemize}

%% Reference
%\newpage
%\bibliographystyle{IEEETran}
%\bibliography{mycite}
\end{document}
